%% Figure 6: Zero-Shot Readiness Experiment
%% KDD 2026 Submission - BanditGPT

\subsection{Zero-Shot Readiness: The "GPT-5" Simulation}
\label{sec:zero-shot-readiness}

To evaluate the system's agility in a production environment where new models are frequently released, we simulated a ``Mid-Flight Upgrade'' event. The router was trained on the standard portfolio (Mixtral-8x7b-Instruct, GPT-4-Turbo) for $t=300$ steps. At $t=300$, we introduced a superior model (``GPT-5.1'', using real held-out evaluation data from the LMSys Arena) into the registry.

\subsubsection{Experimental Design}

We evaluated the system through two complementary experiments:

\paragraph{Ablation Study (Figure~\ref{fig:ablation-study}).} 
We compared three initialization strategies to isolate the contribution of each component:

\begin{itemize}
    \item \textbf{Cold Start (Red):} No warmup priors. All models initialized with $\mathbf{A}=\lambda \mathbf{I}, \mathbf{b}=\vec{0}$. Pure online learning from scratch.
    
    \item \textbf{Warmup Only (Orange):} Existing models (Mixtral, GPT-4-Turbo) initialized with 80k LMSys Arena battle priors. New model (GPT-5.1) added cold at $t=300$ with no transfer.
    
    \item \textbf{Warmup + Semantic Transfer (Green):} Existing models use warmup priors. New model inherits preference vector from semantic neighbor: $\theta_{\text{GPT-5.1}} \leftarrow \theta_{\text{GPT-4-Turbo}}$ while resetting confidence ($\mathbf{A}_{\text{new}} \leftarrow \lambda \mathbf{I}$).
\end{itemize}

\paragraph{Production Router (Figure~\ref{fig:adaptive-efficiency}).}
We evaluate the complete production system with the \textbf{Heterogeneous Experts Strategy}. The router employs two specialized experts managed by a Corralling meta-learner:

\begin{itemize}
    \item \textbf{Conservative Expert (Blue):} Decaying exploration ($\alpha: 1.0 \to 0.01$). Optimizes for exploitation efficiency during stable periods.
    
    \item \textbf{Adaptive Expert (Red):} Constant exploration ($\alpha = 2.0$). Maintains vigilance for distribution shifts and new model releases.
    
    \item \textbf{Meta-Learner:} Corralling algorithm dynamically adjusts expert weights based on cumulative regret, automatically switching between efficiency and exploration modes.
\end{itemize}

\subsubsection{Result Analysis}

\paragraph{Ablation Study Results (Figure~\ref{fig:ablation-study}).} 
The three-way comparison reveals the incremental value of each component:

\textit{Cold Start (Red):} Suffers catastrophic performance regression at $t=300$, plummeting from average reward 3.3 to 1.7. The router explores the new model blindly, leading to poor task assignments. Recovery requires $\sim$500 steps.

\textit{Warmup Only (Orange):} Warmup priors improve burn-in performance during $t=0$-$300$ (faster convergence to optimal model). However, at model release ($t=300$), performance still dips because GPT-5.1 starts cold. Recovery is faster than pure cold start ($\sim$300 steps) because existing models remain strong.

\textit{Warmup + Semantic Transfer (Green):} Exhibits \textbf{zero performance dip} at model release. By inheriting preferences from GPT-4-Turbo, the router immediately exploits GPT-5.1 on relevant tasks. Maintains high reward ($\sim$4.5) from the first timestep after release. Semantic transfer provides a \textbf{$2.8\times$ performance advantage} over cold start during the critical adaptation window ($t=300$ to $t=500$).

\paragraph{Production Router Results (Figure~\ref{fig:adaptive-efficiency}).} 
The heterogeneous experts strategy demonstrates automatic regime adaptation:

\textit{Stable Period ($t=0$-$300$):} Conservative expert (blue) dominates, reaching 85\% weight by $t=300$. Its decaying alpha ($\alpha \to 0.01$) efficiently exploits warmup priors. Adaptive expert (red) remains active but downweighted (15\%), providing vigilance without excessive cost.

\textit{Model Release ($t=300$):} Conservative expert receives transferred knowledge from GPT-4-Turbo via semantic transfer. With $\alpha=0.01$ (pure exploitation), it immediately leverages the transfer for zero-shot GPT-5.1 usage. No performance dip observed.

\textit{Expert Weight Dynamics:} Meta-learner maintains trust in conservative expert post-release (weight remains $\sim$85\%), validating that semantic transfer was successful. Adaptive expert provides exploration backup but is not needed due to effective transfer.

\paragraph{Mechanism of Action.} 
This stability validates our \textbf{Preference-Confidence Decoupling} strategy combined with \textbf{Heterogeneous Exploration}. By transferring the preference vector ($\theta$) but resetting the precision matrix ($\mathbf{A}$), the router starts with a strong hypothesis (``GPT-5.1 likely excels at Math, like GPT-4-Turbo'') but loose confidence intervals (``but I need to verify''). The conservative expert exploits this prior immediately, while the adaptive expert provides a safety net if the transfer fails.

\subsubsection{Practical Implications}

These results have significant implications for production systems:

\begin{enumerate}
    \item \textbf{Cost Savings:} The 500-step recovery period in Cold Start represents wasted API calls and degraded user experience. At scale (millions of queries per day), this translates to substantial cost and quality-of-service penalties. Semantic transfer eliminates this adaptation tax.
    
    \item \textbf{Continuous Deployment:} Zero-shot readiness enables continuous model portfolio updates. Operators can immediately capitalize on new model releases (e.g., GPT-5) without retraining or performance regression.
    
    \item \textbf{Automatic Regime Adaptation:} The heterogeneous experts strategy removes the need for manual hyperparameter tuning. The meta-learner automatically switches between efficiency (stable periods) and exploration (distribution shifts) based on real-time performance.
    
    \item \textbf{Robustness to Model Churn:} As providers deprecate APIs (e.g., GPT-4 $\to$ GPT-4-Turbo $\to$ GPT-4o), semantic transfer ensures seamless transitions. The adaptive expert provides backup vigilance if transferred knowledge degrades.
    
    \item \textbf{Production-Ready Safety:} Unlike homogeneous strategies (uniform alpha decay) that risk ``brain death'' in non-stationary environments, the heterogeneous approach maintains exploration capacity through the adaptive expert while optimizing efficiency through the conservative expert.
\end{enumerate}

%% Figure moved to top of section - see figure6_caption.tex

\subsubsection{Algorithm: Semantic Transfer with Heterogeneous Experts}

Algorithm~\ref{alg:semantic-transfer} formalizes the new model admission procedure used by both experts in the Corralling architecture.

\begin{algorithm}[t]
\caption{Semantic Transfer for New Model Admission}
\label{alg:semantic-transfer}
\begin{algorithmic}[1]
\REQUIRE New model $m_{\text{new}}$, existing portfolio $\mathcal{M}$, embeddings $\{e_m\}_{m \in \mathcal{M}}$
\ENSURE Initialized parameters $(\mathbf{A}_{m_{\text{new}}}, \mathbf{b}_{m_{\text{new}}})$

\STATE Compute embedding: $e_{\text{new}} \leftarrow \text{SentenceTransformer}(\text{description}(m_{\text{new}}))$

\STATE Find semantic neighbor:
\[
m^* = \arg\max_{m \in \mathcal{M}} \text{cosine\_similarity}(e_{\text{new}}, e_m)
\]

\STATE Extract learned preference:
\[
\theta^* = \mathbf{A}_{m^*}^{-1} \mathbf{b}_{m^*}
\]

\STATE \textbf{Initialize with transfer:}
\STATE \quad $\mathbf{A}_{m_{\text{new}}} \leftarrow \lambda \mathbf{I}$ \COMMENT{Reset confidence (high exploration)}
\STATE \quad $\mathbf{b}_{m_{\text{new}}} \leftarrow N_{\text{eff}} \cdot \theta^*$ \COMMENT{Transfer intuition ($N_{\text{eff}} \approx 5$)}

\RETURN $(\mathbf{A}_{m_{\text{new}}}, \mathbf{b}_{m_{\text{new}}})$
\end{algorithmic}
\end{algorithm}

\paragraph{Key Parameters.}
We set $N_{\text{eff}} = 5.0$, reflecting the belief that the semantic neighbor provides information equivalent to 5 observed samples. This hyperparameter controls the strength of the prior: higher values increase initial exploitation, while lower values allow faster adaptation if the neighbor differs significantly from the new model.

\subsubsection{Theoretical Justification}

The efficacy of our approach rests on three key principles:

\begin{enumerate}
    \item \textbf{Task-Capability Correlation:} Models with similar architectures/training exhibit correlated performance across task types. For example, GPT-4-Turbo and GPT-5.1 both excel at reasoning tasks (Math, Code) but may struggle with creative writing compared to Claude-3.5-Sonnet. This justifies semantic transfer.
    
    \item \textbf{Semantic Embedding Validity:} SentenceTransformer embeddings of model descriptions capture meaningful capability distinctions. Embedding-based neighbor selection outperforms random selection by $37\%$ in post-release regret.
    
    \item \textbf{Heterogeneous Exploration Principle:} In non-stationary environments, no single exploration schedule is universally optimal. The Corralling meta-learner (Agarwal et al., 2017) provides sublinear regret guarantees by adaptively combining expert policies. Our heterogeneous strategy instantiates this framework with complementary exploration philosophies.
\end{enumerate}

The transferred preference vector $\theta^*$ encodes the \textit{latent task affinity} learned from the neighbor, providing a warm start that is calibrated to the task distribution. The conservative expert ($\alpha \to 0.01$) immediately exploits this knowledge, while the adaptive expert ($\alpha = 2.0$) provides a corrective mechanism if the transfer degrades. The meta-learner automatically reweights experts based on observed regret, ensuring robustness without manual intervention.

